% BHU PYQ -- Manually transcribed & error-corrected
% Paper: BCH-122 : Financial Accounting-II
% Exam: B.Com. (Hons.) Semester II Examination, 2015-16

\documentclass[11pt,a4paper]{article}
\usepackage[a4paper,top=2.5cm,bottom=2.5cm,left=2.8cm,right=2.8cm,headheight=14pt]{geometry}
\usepackage{amsmath,amssymb}
\usepackage[T1]{fontenc}
\usepackage[utf8]{inputenc}
\usepackage{lmodern,microtype,fancyhdr,enumitem,hyperref,lastpage,parskip}
\usepackage{booktabs}

\hypersetup{
    colorlinks=true,
    urlcolor=black,
    linkcolor=black,
    pdftitle={BCH-122: Financial Accounting-II -- BHU B.Com. (Hons.) Sem II 2015-16}
}

\pagestyle{fancy}
\fancyhf{}
\renewcommand{\headrulewidth}{0.4pt}
\renewcommand{\footrulewidth}{0.4pt}
\fancyhead[L]{\small\textit{Banaras Hindu University}}
\fancyhead[R]{\small\textit{BCH-122: Financial Accounting-II}}
\fancyfoot[C]{\small Page \thepage\ of \pageref{LastPage}}

\newlist{parts}{enumerate}{2}
\setlist[parts,1]{label=(\alph*),leftmargin=*,itemsep=5pt,topsep=3pt}
\setlist[parts,2]{label=(\roman*),leftmargin=*,itemsep=3pt,topsep=2pt}
\newcommand{\pts}[1]{\hfill\textit{[#1]}}

\begin{document}

\begin{center}
  {\large\bfseries BANARAS HINDU UNIVERSITY}\\[2pt]
  {\normalsize B.Com. (Hons.) --- Semester II Examination, 2015-16}\\[6pt]
  \rule{0.8\linewidth}{0.5pt}\\[6pt]
  {\large\bfseries Paper No. BCH-122: Financial Accounting-II}\\[4pt]
  {\normalsize Time: 3 Hours\hfill Maximum Marks: 70}
\end{center}

\rule{\linewidth}{0.4pt}

\smallskip
\noindent\textit{Instructions: Answer all questions. All questions carry equal marks (14 marks each).}

\bigskip

\noindent\textbf{Question 1.}
\begin{parts}
  \item Define Bill of Exchange. Explain the various types of Bill of Exchange.\pts{5}
  \item On 1st July, 2015, Mahesh sold goods to Naresh for \pounds 25,000 (read as Rupees). Naresh paid \pounds 5,000 in cash and accepted a Bill of Exchange for the balance of \pounds 20,000 for a period of three months. Mahesh discounted the bill with his banker at 10\% p.a. On the due date, Naresh failed to pay the bill and the bank paid \pounds 100 for noting charges. Subsequently, Naresh offered to pay \pounds 10,000 against the dishonoured bill and also the noting charges in cash, and to accept a new bill for the balance amount of \pounds 10,000 for a period of two months. Mahesh agreed to the request. The new bill was discharged by Naresh on the due date. Journalise the above transactions in the books of both the parties.\pts{9}
\end{parts}

\centerline{\textbf{OR}}

What is a Bank Reconciliation Statement? Explain its utility and the reasons for discrepancies between the cash book and pass book balances.\pts{14}

\medskip\rule{\linewidth}{0.2pt}

\newpage

\noindent\textbf{Question 2.}
X Ltd. was registered with an authorized capital of \pounds 20,00,000 divided into 20,000 shares of \pounds 100 each. 10,000 shares were offered to the public at a premium of \pounds 20 per share, payable as follows:
\begin{itemize}[itemsep=2pt,topsep=2pt]
  \item On Application: \pounds 20
  \item On Allotment: \pounds 45 (including premium)
  \item On First Call: \pounds 25
  \item On Second and Final Call: Balance amount
\end{itemize}
Applications were received for 14,000 shares. Pro-rata allotment was made to the applicants of 12,000 shares, and the remaining applications were rejected and their application money was returned. Prakash, holding 200 shares, paid the entire call money along with allotment. Pradeep, holding 300 shares, failed to pay the allotment and first call money, and his shares were subsequently forfeited. These shares were reissued as fully paid-up at \pounds 90 per share. Pass the necessary journal entries in the books of the company.\pts{14}

\centerline{\textbf{OR}}

\begin{parts}
  \item Explain the legal provisions regarding the redemption of preference shares.\pts{5}
  \item Y Ltd. has issued 4,00,000; 14\% redeemable preference shares of \pounds 100 each. Redemption was to take place on 1 April, 2016. A general reserve of \pounds 2,50,00,000 has already been built up out of past profits. For the purpose of redemption, 75,000 new 15\% preference shares of \pounds 200 each at a premium of \pounds 100 each, were offered to the public in full. The new issue was subscribed and paid for. Thereafter, the 14\% redeemable preference shares were redeemed. Pass the journal entries.\pts{9}
\end{parts}

\medskip\rule{\linewidth}{0.2pt}

\noindent\textbf{Question 3.}
Distinguish between a share and a debenture. Describe the various methods of redemption of debentures.\pts{14}

\centerline{\textbf{OR}}

On 1st January, 2016, Genuine Products Ltd. issued 1,000; 8\% debentures of \pounds 100 each, which are redeemable after four years at par. For the purpose of redemption, a non-cumulative 8\% debentures sinking fund is to be created. The sinking fund investments were sold for \pounds 78,000 on December 31, 2019. On this date, the bank balance stood at \pounds 42,000 (before the sale of investments). The interest received on sinking fund investments for the years 2017, 2018, and 2019 were \pounds 2,000, \pounds 4,000, and \pounds 6,000. Prepare the Sinking Fund A/c, Sinking Fund Investment A/c, Bank A/c, and Interest on Sinking Fund Investment A/c.\pts{14}

\medskip\rule{\linewidth}{0.2pt}

\newpage

\noindent\textbf{Question 4.}
XYZ Ltd. was formed to take over a running business from 1st April, 2015. The company was incorporated on 1st August, 2015. At the end of the year on 31st March, 2016, the company prepared the following Profit and Loss Account:

\begin{center}
\begin{tabular}{lr|lr}
\toprule
\textbf{Particulars} & \textbf{Amount (\pounds)} & \textbf{Particulars} & \textbf{Amount (\pounds)} \\
\midrule
To Rent & 30,000 & By Gross Profit b/d & 3,36,000 \\
To Insurance & 10,000 & & \\
To Electricity & 4,800 & & \\
To Salaries & 72,000 & & \\
To Director Fees & 9,200 & & \\
To Commission & 12,000 & & \\
To Advertisement & 8,000 & & \\
To Office Expenses & 15,000 & & \\
To Depreciation & 6,000 & & \\
To Preliminary Expenses & 15,000 & & \\
To Bad Debts & 4,000 & & \\
To Interest on Loan & 6,000 & & \\
To Net Profit & 1,44,000 & & \\
\midrule
\textbf{Total} & \textbf{3,36,000} & \textbf{Total} & \textbf{3,36,000} \\
\bottomrule
\end{tabular}
\end{center}

\noindent\textbf{Additional Information:}
\begin{parts}
  \item Sales up to the incorporation of the business amounted to \pounds 4,80,000 whereas they spurted to record an increase of two-fifths during the rest of the year.
  \item Rent included the rent of the building paid @ \pounds 2,000 per month up to September, 2015.
  \item Depreciation included \pounds 1,200 for the assets acquired in the post-incorporation period.
  \item The purchase consideration was discharged by the company on 30th September, 2015 by issuing equity shares of \pounds 10 each.
\end{parts}
Prepare the Profit and Loss Account showing the allocation of profit between the pre-incorporation and post-incorporation periods, indicating the basis of allocation.\pts{14}

\centerline{\textbf{OR}}

Draft the Balance Sheet of a company in both Horizontal and Vertical forms as per the Companies Act with imaginary figures.\pts{14}

\medskip\rule{\linewidth}{0.2pt}

\newpage

\noindent\textbf{Question 5.}
Following is the Trial Balance of Mysore Ltd. as on 31st March, 2016:

\begin{center}
\begin{tabular}{lr|lr}
\toprule
\textbf{Debit Balances} & \textbf{Amount (\pounds)} & \textbf{Credit Balances} & \textbf{Amount (\pounds)} \\
\midrule
Land \& Building & 34,000 & Share Capital & 80,000 \\
Plant \& Machinery & 66,000 & Reserve Fund & 12,000 \\
Fees to Credit Rating Agency & 2,000 & Pension Fund & 6,000 \\
Calls in Arrear & 600 & Profit \& Loss A/c & 3,540 \\
5\% Govt. Bonds & 3,600 & Return Outward & 2,000 \\
Motor Vehicles & 4,000 & Sales & 1,23,000 \\
Interim Dividend & 1,800 & Suppliers of Materials & 10,000 \\
Goodwill & 54,600 & Debentures & 30,000 \\
Repairs & 300 & Income Tax Outstanding & 5,000 \\
Bank Interest & 840 & Foreign Exchange Fluctuation Fund & 45,000 \\
Purchases & 96,000 & & \\
Audit Fees & 400 & & \\
Advertisement & 1,000 & & \\
Carriage Inward & 1,500 & & \\
Wages & 9,200 & & \\
Insurance & 2,000 & & \\
Opening Inventory & 19,000 & & \\
Tools & 4,000 & & \\
Loss on Foreign Exchange & 1,700 & & \\
Furniture & 2,900 & & \\
Trade Debtors & 8,300 & & \\
Return Inward & 2,800 & & \\
\midrule
\textbf{Total} & \textbf{3,16,540} & \textbf{Total} & \textbf{3,16,540} \\
\bottomrule
\end{tabular}
\end{center}

\noindent\textbf{Adjustments:}
\begin{parts}
  \item Closing inventory was valued at \pounds 20,000.
  \item Company spent \pounds 5,600 for protection of ecology and factory environment, which has not yet been accounted for.
  \item Manager's commission to be paid @ 2\% of Gross Profit before charging such commission.
  \item Outstanding liabilities: Wages \pounds 300.
  \item Write off \pounds 2,000 from Goodwill.
  \item Printing expenses for the previous year of \pounds 300 were not provided in that year, now paid in the current year.
\end{parts}
Prepare a Trading and Profit and Loss Account for the year ended 31st March, 2016 and also a Balance Sheet as on that date.\pts{14}

\centerline{\textbf{OR}}

\begin{parts}
  \item Explain the provisions of the Companies Act regarding the preparation and presentation of the Profit and Loss Account.\pts{9}
  \item Give the proforma of the Profit and Loss A/c of a company as per the provisions of the Companies Act.\pts{5}
\end{parts}

\vfill
\rule{\linewidth}{0.4pt}
\begin{center}\small
\href{https://bhu-exam.vercel.app/}{Click here to visit our website}\\[4pt]
\href{https://me-aryan.vercel.app/}{Click here to visit our contributor}\\[4pt]
\href{https://quashan-me.vercel.app/}{Click here to visit our contributor}
\end{center}

\end{document}
